\documentclass[openright, oneside, a4paper, article, brazil]{abntex2}
\usepackage{formularios}
% --------- %
% Metadados %
% --------- %
\hypersetup{%
    colorlinks=false,
    allbordercolors=white,
    pdftitle = {Requerimento de inscrição de pessoa jurídica},
    pdfauthor = {Conselho Regional de Psicologia de São Paulo},
    pdfsubject = {inscrição},
    pdfkeywords = {%
            inscrição,
            registro,
            pessoa jurídica,
            PJ,
        }%
}
\title{Requerimento de inscrição de pessoa jurídica}
\author{Conselho Regional de Psicologia de São Paulo}
% --------- %
% DOCUMENTO %
% --------- %
\begin{document}
\pagestyle{crp}
\centering
\Form
\section*{REQUERIMENTO DE INSCRIÇÃO DE PESSOA JURÍDICA}%
% -------------------- %
% EXEMPLOS DE CONTEÚDO %
% -------------------- %
% -------------------------- %
% Identificação/qualificação %
% -------------------------- %

\setasuspacing{\DoubleSpacing}
% Define espaçamento duplo para que as caixas de texto não fiquem sobrepostas
\justifying
% Parágrafo justificado
\textbf{À/ao presidenta/e do Conselho Regional de Psicologia da 6ª Região (CRP-06)}
\vspace{\baselineskip}

\noindent A pessoa jurídica \preenctext{14}{razao_social}{34.3em}{(razão social) } \hfill,\\
\preenctext{14}{nome_fantasia}{20em}{de nome de fantasia} \hfill, \hfill \preenctext{14}{cnpj}{15em}{(CNPJ) } \hfill,\\
\preenctext{14}{cnpj}{15em}{CNPJ nº } , \preenctext{14}{endereco}{15em}{com sede à} \hfill,\\
\preenctext{14}{numero}{4em}{nº }, \preenctext{14}{bairro}{15em}{bairro } , \preenctext{14}{cep}{8em}{CEP }, \preenctext{14}{cidade}{15em}{cidade }	 \hfill,\\
\preenctext{14}{uf}{2em}{UF } , \preenctext{14}{telefone}{15em}{telefone } , \preenctext{14}{ramal}{3em}{ramal } , \preenctext{14}{email}{15em}{e-mail } \hfill,\\
\preenctext{14}{site}{12em}{site } , por seu \preenctext{14}{representante}{15em}{representante legal } \hfill,\\
\preenctext{14}{cpf}{15em}{CPF} , \preenctext{14}{tel_rep}{15em}{telefone } ,\hfill,\\
\noindent
abaixo assinada(o), conforme o disposto na Lei nº 5.766, de 20 de dezembro de 1971 e no Decreto nº 79.822,
de 17 de junho de 1977, que a regulamenta, e na Consolidação das Resoluções do Conselho Federal de Psicologia,
vem respeitosamente à presença de V. Sa. requerer a INSCRIÇÃO DE PESSOA JURÍDICA, por ter como objetivo
social a prestação de serviços psicológicos a terceiros ou por ter psicóloga/o na equipe de trabalho.

Para tal requerimento, declaro ter anexado os seguintes documentos obrigatórios:

\begin{enumerate}
    \item \textbf{instrumento de constituição}	(contrato social, requerimento de empresário, ata ou estatutos)
    da pessoa jurídica consolidado com as últimas alterações, regis­tradas em cartório competente ou na
    Junta Comercial (será aceita autenticação digital dos documen­tos da Junta Comercial);
    \item \textbf{documento que atribua poderes ao representante legal} (ex: ata de eleição e posse da
    diretoria, contrato social, procuração);
    \item \textbf{requerimento de inscrição/ficha cadastral de pessoa jurídica} devidamente preenchida/o
    e assinada/o pela/o psicóloga/o indicada/o como Responsável Técnica/o e a/o representante legal
    da pessoa jurídica (será aceita assinatura com certificação digital. Ex: gov.br);
    \item Termo de Responsabilidade Técnica e Autorização para Envio de Comunicações Eletrônicas,
    devidamente preenchido e assinado pela(o) Psicóloga(o) indicada(o) para a função (Será aceita assinatura
    com certificação digital. Ex: Gov.br);
    \item Declaração Institucional para Garantia do Amplo e Livre Exercício Profissional às/os
    psicólogas/os e responsáveis técnicas/os de psicologia, de acordo com o Código de Ética profissional
    da/o psicóloga/o e demais legislações (Será aceita assinatura com certificação digital. Ex: Gov.br);
    \item Comprovante de vínculo de trabalho da (o)(s) responsável(eis) técnica (o)(s), por meio dos seguintes
    documentos (carteira de trabalho, contrato de prestação de serviços, termo de adesão ao trabalho voluntário
    ou documento constitutivo da empresa, quando sócio);
    \item Se houver estagiários de Psicologia, cópia dos termos de compromisso firmados com as Universida­des ou Faculdades;
    \item Cartão CNPJ ou documento emitido pela Internet (impressão atualizada do comprovante de inscrição CNPJ);
    \item Se a entidade for filantrópica, documento que comprove ser de utilidade públi­ca, estatuto e outros,
    devidamente registrados em Cartório, se houver;
    \item Documento que regulamente as normas de funcionamento (ex: regimento interno), se houver;
    \item Certificado de Registro da Pessoa Jurídica em outro Conselho de fiscalização Profissional, se houver;
    \item Certificado de filantropia emitido pelo CNAS (Conselho Nacional de Assistência Social), se houver;
    \item Alteração Contratual que tenha ocorrido desde a constituição legal da Pessoa Jurídica até a presente
    data, autenticada;
    \item Alvará ou Certificado de Licenciamento Integrado (CLI).
\end{enumerate}
\vspace{4\baselineskip}

\centering
\dataextenso{1}
\vspace{3\baselineskip}

\centering
\rule{24em}{0.4pt}\\
\small Assinatura da/do representante legal da pessoa jurídica
\vfill{}

\end{document}
